\section{Discussion}
Evaluation mattered more than architecture here. Ensembling, fusion, heavy augmentation, extra seeds and pseudo-labelling each moved the held-out score by about 0.01 at most, none of them significantly, and seed noise was about as large. Papers that report differences of this size without intervals deserve some caution.
Most errors sit in the rare classes (1 to 6 held-out positives) and in the heterogeneous "Others" class. The Hip-to-Pelvis cluster in the audit, and the competition hosts' offline re-annotation of the test set, suggest that label noise sets part of the ceiling, as Langer et al. (2024) found on a hospital dataset. The label-efficiency curve was still rising, so more labelled data would probably help more than a different architecture.
We did not test why the gate collapsed. One possibility is that a gate learned on the raw features of two networks with different activation scales lets one branch dominate. The learned weights point to a degenerate solution, not to complementary fusion, and fixed averaging was no worse.
For clinical use, a router that answers 90 \% of the images keeps 94.5 \% top-1 accuracy and sends the other 10 \% to a person. The operating point chosen on OOF predictions should be re-validated on local data before anyone relies on it. We see the classifier as a way to flag suspicious metadata, not as a replacement for it.
Our MURA accuracy is higher than the one published figure we found for this task. Other studies report 97 to 99 \% on other data and with much larger private training sets (Langer et al. 2024; Mitsuyama et al. 2026), and we do not claim to outperform them.

